\section*{अध्याय \ref{ch:g12:deriv} --- अवकलन और उत्तलता}

\begin{solution}{exo:g12:deriv:1}
$f'(x) = 3x^2 - 5$।

भागफल का नियम: $g'(x) = \dfrac{(x^2+1) - x\cdot 2x}{(x^2+1)^2} = \dfrac{1 - x^2}{(x^2+1)^2}$।

$u = 2x+1$ के साथ शृंखला नियम: $h'(x) = 5 \cdot 2 \cdot (2x+1)^4 = 10(2x+1)^4$।

$u = x^2 + 1$ के साथ शृंखला नियम: $k'(x) = \dfrac{2x}{2\sqrt{x^2+1}} = \dfrac{x}{\sqrt{x^2+1}}$।
\end{solution}

\begin{solution}{exo:g12:deriv:2}
$f'(x) = 2x - 3$, इसलिए $f(2) = -1$ और $f'(2) = 1$: \omterm{def:g12:deriv:tangent}{स्पर्श रेखा} $y = -1 + (x - 2) = x - 3$ है। अंतर
\[
f(x) - (x - 3) = x^2 - 4x + 4 = (x-2)^2 \geq 0 ,
\]
है, इसलिए वक्र हर जगह \omterm{def:g12:deriv:tangent}{स्पर्श रेखा} के ऊपर पड़ता है और उसे केवल $x = 2$ पर छूता है
(जैसी अपेक्षा थी: $f$ \omterm{def:g12:deriv:convex}{उत्तल} है)।
\end{solution}

\begin{solution}{exo:g12:deriv:3}
दोनों अंतर-भागफल हैं। $0$ पर $f(x) = (1+x)^{100}$ के साथ: $f'(x) = 100(1+x)^{99}$, इसलिए सीमा $f'(0) = 100$ है। $4$ पर
$g(x) = \sqrt{x}$ के साथ: $g'(x) = \frac{1}{2\sqrt{x}}$, इसलिए सीमा $g'(4) = \frac14$ है।
\end{solution}

\begin{solution}{exo:g12:deriv:4}
\omterm{def:g11:func:function}{प्रांत} $\R \setminus \{1\}$। सीमाएँ: $\pm\infty$ पर $f(x) \sim x \to
\pm\infty$; $1^\pm$ पर अंश $\to 4 > 0$ और हर $\to 0^\pm$, इसलिए $f \to \pm\infty$;
रेखा $x = 1$ \omterm{def:g12:limcont:asymptote}{ऊर्ध्वाधर अनन्तस्पर्शी} है।

\[
f'(x) = \frac{2x(x-1) - (x^2+3)}{(x-1)^2} = \frac{x^2 - 2x - 3}{(x-1)^2}
= \frac{(x-3)(x+1)}{(x-1)^2}.
\]
अतः $\intoo{-\infty}{-1}$ और $\intoo{3}{+\infty}$ पर $f' > 0$, तथा $\intoo{-1}{1}$ और $\intoo{1}{3}$ पर $f' < 0$। \omterm{def:g11:func:function}{फलन} बढ़कर \omterm{def:g12:deriv:extremum}{स्थानीय अधिकतम}
$f(-1) = -2$ तक जाता है, घटकर $-\infty$ तक आता है, \omterm{def:g12:limcont:asymptote}{अनन्तस्पर्शी} के बाद उछलकर $+\infty$ पर आता
है, घटकर स्थानीय \omterm{def:g10:functions:extrema}{न्यूनतम} $f(3) = 6$ तक जाता है, और फिर बढ़ता है।
\end{solution}

\begin{solution}{exo:g12:deriv:5}
$x \in \intoo{0}{15}$ के लिए डिब्बे का आधार भुजा $30 - 2x$ का वर्ग है और ऊँचाई $x$, इसलिए
\[
V(x) = x(30 - 2x)^2 .
\]
$V'(x) = (30-2x)^2 + x \cdot 2(30-2x)(-2) = (30-2x)(30 - 2x - 4x)
= (30-2x)(30-6x)$। $\intoo{0}{15}$ पर $30 - 2x > 0$, इसलिए $V'$ का चिह्न $30 - 6x$ जैसा है: $x = 5$ से पहले धनात्मक,
बाद में ऋणात्मक। आयतन $x = 5$ cm पर \omterm{def:g10:functions:extrema}{अधिकतम} होता है, और वह $V(5) = 5 \times 20^2 = 2000\ \text{cm}^3$ है।
\end{solution}

\begin{solution}{exo:g12:deriv:6}
\emph{1.} $f'(x) = 4x^3 - 12x^2$ और $f''(x) = 12x^2 - 24x = 12x(x-2)$। इस प्रकार $\intoo{-\infty}{0}$ और $\intoo{2}{+\infty}$ पर $f'' > 0$ (\omterm{def:g12:deriv:convex}{उत्तल}), तथा $\intoo{0}{2}$ पर
$f'' < 0$ (\omterm{def:g12:deriv:convex}{अवतल})। $f''$ $0$ पर और $2$ पर चिह्न बदलता है: दोनों \omterm{def:g12:deriv:inflection}{नति-परिवर्तन बिंदु}
हैं।

\emph{2.} $a = 2$ पर: $f(2) = 16 - 32 + 10 = -6$, $f'(2) = 32 - 48 = -16$, \omterm{def:g12:deriv:tangent}{स्पर्श रेखा} $y = -6 - 16(x-2) = -16x + 26$। अंतर
\[
g(x) = f(x) + 16x - 26 = x^4 - 4x^3 + 16x - 16 .
\]
है। चूँकि $g(2) = g'(2) = g''(2) = 0$ है, इसलिए $(x-2)^3$ $g$ को विभाजित करता है; भाग देने पर $g(x) = (x-2)^3(x+2)$ मिलता
है। $x = 2$ के पास $x + 2 > 0$, इसलिए $g$ का चिह्न $(x-2)^3$ जैसा है: $2$ से पहले ऋणात्मक,
बाद में धनात्मक। वक्र \omterm{def:g12:deriv:tangent}{स्पर्श रेखा} को पार कर जाता है: $2$ सचमुच
\omterm{def:g12:deriv:inflection}{नति-परिवर्तन बिंदु} है।
\end{solution}

\begin{solution}{exo:g12:deriv:7}
\omterm{def:g11:func:function}{फलन} $f(x) = \sqrt{x}$ $\intoo{0}{+\infty}$ पर $f''(x) = -\frac{1}{4}x^{-3/2} < 0$ संतुष्ट करता है: वह \omterm{def:g12:deriv:convex}{अवतल} है, इसलिए उसका \omterm{def:g10:functions:graph}{आलेख} अपनी हर
\omterm{def:g12:deriv:tangent}{स्पर्श रेखा} के \emph{नीचे} पड़ता है। $a = 1$ पर \omterm{def:g12:deriv:tangent}{स्पर्श रेखा}
\[
y = f(1) + f'(1)(x - 1) = 1 + \tfrac12 (x-1) = \frac{x+1}{2},
\]
है, जिससे सभी $x > 0$ के लिए $\sqrt{x} \leq \frac{x+1}{2}$ निकलता है। समता का अर्थ है कि वक्र \omterm{def:g12:deriv:tangent}{स्पर्श रेखा}
को छूता है, और यह केवल स्पर्श बिंदु $x = 1$ पर होता है। (समतुल्य रूप से: $\left(\sqrt x - 1\right)^2 \geq 0$।)
\end{solution}

\begin{solution}{exo:g12:deriv:8}
\cref{thm:g12:deriv:convexity} से किसी \omterm{def:g12:deriv:derivative}{अवकलनीय} \omterm{def:g12:deriv:convex}{उत्तल} \omterm{def:g11:func:function}{फलन} का \omterm{def:g10:functions:graph}{आलेख} अपनी हर \omterm{def:g12:deriv:tangent}{स्पर्श रेखा} के ऊपर पड़ता है।
$a$ पर \omterm{def:g12:deriv:tangent}{स्पर्श रेखा} क्षैतिज है ($f'(a) = 0$), जिसका \omterm{def:g10:algebra:equation}{समीकरण} $y = f(a)$ है। अतः सभी $x \in \R$
के लिए $f(x) \geq f(a)$: $f(a)$ वैश्विक \omterm{def:g10:functions:extrema}{न्यूनतम} है।
\end{solution}

\begin{solution}{exo:g12:deriv:9}
\emph{नियत परिमाप:} भुजाओं $x$ और $\frac{p}{2} - x$ ($0 < x < \frac p2$) वाले आयत का क्षेत्रफल $S(x) = x\left(\frac p2 - x\right)$
है। तब $S'(x) = \frac p2 - 2x$ $x = \frac p4$ पर लुप्त होता है, और उससे पहले धनात्मक तथा बाद में ऋणात्मक
है: \omterm{def:g10:functions:extrema}{अधिकतम} $x = \frac{p}{4}$ पर है, जहाँ दोनों भुजाएँ $\frac p4$ के बराबर हैं --- यानी वर्ग।

\emph{नियत क्षेत्रफल:} भुजाएँ $x$ और $\frac Ax$ परिमाप $P(x) = 2\left(x + \frac Ax\right)$ देती हैं, जिसमें
$P'(x) = 2\left(1 - \frac{A}{x^2}\right)$ $x < \sqrt A$ के लिए ऋणात्मक और उसके बाद धनात्मक है: \omterm{def:g10:functions:extrema}{न्यूनतम} $x = \sqrt{A}$ पर, फिर एक
वर्ग।

\emph{संबंध:} दोनों कथन द्वैत हैं। मान लीजिए नियत क्षेत्रफल $A$ पर परिमाप
को कोई अ-वर्ग आयत \omterm{def:g10:functions:extrema}{न्यूनतम} करता; तब पहले कथन से उसी परिमाप वाले वर्ग का
क्षेत्रफल $> A$ होता, और उसे समरूप रूप से सिकोड़कर क्षेत्रफल $A$ तक लाने
पर उसका परिमाप कड़ाई से घट जाता --- जो न्यूनतमता का खंडन है। इस प्रकार हर कथन
दूसरे को सिद्ध कर देता है।
\end{solution}

\begin{solution}{pb:g12:deriv:1}
\textbf{1.} $30x\left(3x^2 + 1\right)^4$; \quad $\dfrac{x}{\sqrt{x^2 + 9}}$; \quad $-\dfrac{2x}{\left(x^2 + 1\right)^2}$।

\textbf{2.} $y(4) = 5$, \omterm{def:g10:reffunc:affine}{प्रवणता} $\frac45$: \omterm{def:g12:deriv:tangent}{स्पर्श रेखा} $y = 5 + \frac45(x - 4)$, अर्थात् $y = \frac45 x + \frac95$।

\textbf{3.} $f'(x) = \frac{(x^2 + 1) - x \cdot 2x}{(x^2+1)^2}
= \frac{1 - x^2}{(x^2 + 1)^2}$: $-1$ और $1$ के आरपार ऋणात्मक, धनात्मक, ऋणात्मक:
\omterm{def:g10:functions:extrema}{न्यूनतम} $f(-1) = -\frac12$, \omterm{def:g10:functions:extrema}{अधिकतम} $f(1) = \frac12$, और दोनों ओर \omterm{def:g12:limcont:asymptote}{क्षैतिज अनन्तस्पर्शी} $0$।

\textbf{4.} $g'(x) = 3x^2 - 6x = 3x(x - 2)$: $0$ और $2$ के आरपार \omterm{def:g12:seq:monotonic}{वर्धमान}, \omterm{def:g10:functions:variations}{ह्रासमान}, \omterm{def:g12:seq:monotonic}{वर्धमान};
\omterm{def:g12:deriv:extremum}{स्थानीय अधिकतम} $g(0) = 4$, स्थानीय \omterm{def:g10:functions:extrema}{न्यूनतम} $g(2) = 0$। $g''(x) = 6x - 6$: $1$ से पहले \omterm{def:g12:deriv:convex}{अवतल}, बाद
में \omterm{def:g12:deriv:convex}{उत्तल}: $(1, 2)$ पर नति-परिवर्तन।

\textbf{5.} $1$ पर \omterm{def:g12:deriv:tangent}{स्पर्श रेखा}: \omterm{def:g10:reffunc:affine}{प्रवणता} $g'(1) = -3$: $y = 2 - 3(x - 1) = -3x + 5$। अंतर: $x^3 - 3x^2 + 4 - (-3x + 5) = x^3 - 3x^2 + 3x - 1 =
(x - 1)^3$, जो $1$
पर चिह्न बदलता है: \omterm{def:g12:deriv:tangent}{स्पर्श रेखा} वक्र को \emph{पार} कर जाती है --- यही
\omterm{def:g12:deriv:inflection}{नति-परिवर्तन बिंदु} का विशिष्ट व्यवहार है, जहाँ वक्र पाला बदल लेता है।

\textbf{6.} $T(x) = \dfrac{\sqrt{x^2 + 900}}{5} +
\dfrac{\sqrt{(40 - x)^2 + 400}}{2}$। $0$ से पहले या $40$ से आगे पानी में उतरने से \emph{दोनों}
चरण लंबे हो जाते हैं: व्यर्थ।

\textbf{7.} $T'(x) = \dfrac{x}{5\sqrt{x^2 + 900}} -
\dfrac{40 - x}{2\sqrt{(40 - x)^2 + 400}}$।

\textbf{8.} दौड़ वाले त्रिभुज में तट के अनुदिश भुजा $x$ है और कर्ण $\sqrt{x^2 + 900}$:
लंब के साथ बने कोण की \omterm{def:g11:trigo:cossin}{ज्या} ठीक इन का भागफल है (सामने वाली भुजा बटा कर्ण);
इसी तरह $\sin\theta_2 = \frac{40 - x}{\sqrt{(40-x)^2 + 400}}$। तब $T'(x) = 0$ का अर्थ $\frac{\sin\theta_1}{5} = \frac{\sin\theta_2}{2}$ हो जाता है --- यानी चालों $5$ और
$2$ के साथ स्नेल का नियम।

\textbf{9.} $T(24) \approx 20.5$ s; $T(40) = 10 + 10 = 20$ s; $T(0) = 6 + \sqrt{2000}/2 \approx 28.4$ s। सीधी रेखा ``पूरे रास्ते दौड़िए, फिर
सीधे तैरिए'' से भी हार जाती है --- और दोनों अपवर्तित रास्ते से हार जाते हैं।

\textbf{10.} $T'$ पर \omterm{thm:g12:limcont:ivt}{द्विभाजन} ($24$ पर ऋणात्मक, $40$ पर धनात्मक) $x^* \approx 34$ m
पर \omterm{def:g12:seq:limit}{अभिसरित} होता है, जिसमें $T(x^*) \approx 19.5$ s है: सीधी रेखा से लगभग एक सेकंड तेज़ ---
और यही एक सेकंड बचाव और त्रासदी के बीच का अंतर है।

\textbf{11.} हर चरण का समय $x$ का \omterm{def:g12:deriv:convex}{उत्तल} \omterm{def:g11:func:function}{फलन} है (एक $\sqrt{\text{द्विघात}}$ शाखा), इसलिए $T$
\omterm{def:g12:deriv:convex}{उत्तल} है, और किसी \omterm{def:g12:deriv:derivative}{अवकलनीय} \omterm{def:g12:deriv:convex}{उत्तल} \omterm{def:g11:func:function}{फलन} का क्रांतिक बिंदु उसका वैश्विक \omterm{def:g10:functions:extrema}{न्यूनतम}
होता है (\cref{exo:g12:deriv:8}): $x^*$ केवल स्थिर नहीं, अजेय है।

\textbf{12.} पानी में प्रकाश धीमा है; फ़र्मा के सिद्धांत से हवा-से-पानी वाला
सबसे तेज़ रास्ता सतह पर ठीक स्नेल के नियम के अनुसार मुड़ जाता है --- इसलिए
डूबी हुई पेंसिल से आने वाली किरणें आँख तक मुड़कर पहुँचती हैं, और मस्तिष्क,
सीधी रेखाएँ आगे बढ़ाकर, पेंसिल को जल-रेखा पर टूटा हुआ देख लेता है।

\textbf{13.} $(x^2)'' = 2 > 0$: \omterm{def:g12:deriv:convex}{उत्तल}। \omterm{prop:g10:coordgeom:midpoint}{मध्यबिंदु} पर \omterm{def:g12:deriv:convex}{उत्तलता}: $f\left(\frac{a+b}{2}\right) \leq
\frac{f(a) + f(b)}{2}$, और यही दावा है।
हाथ से: $\frac{a^2 + b^2}{2} - \left(\frac{a+b}{2}\right)^2 =
\frac{(a - b)^2}{4} \geq 0$।

\textbf{14.} प्रश्न 13 में वर्गमूल (\omterm{def:g12:seq:monotonic}{वर्धमान}) लेने पर: $\frac{a + b}{2} \leq \sqrt{\frac{a^2 + b^2}{2}}$। $2, 8$ पर पूरी
कड़ी: हरात्मक $\frac{2 \cdot 16}{10} = 3.2$; \omterm{def:g12:seq:arith-geom}{गुणोत्तर} $\sqrt{16} = 4$; \omterm{def:g12:seq:arith-geom}{समांतर} $5$; वर्ग $\sqrt{34} \approx 5.83$: हर \omterm{def:g11:stat:mean}{माध्य} अगले
के आगे झुक जाता है।

\textbf{15.} $1$ पर $\frac1x$ की \omterm{def:g12:deriv:tangent}{स्पर्श रेखा}: मान $1$, \omterm{def:g10:reffunc:affine}{प्रवणता} $-1$: $y = 2 - x$।
\omterm{def:g12:deriv:convex}{उत्तल} वक्र अपनी स्पर्श रेखाओं के ऊपर बैठते हैं (\cref{thm:g12:deriv:convexity}): सभी $x > 0$ के लिए
$\frac1x \geq 2 - x$, और समता ठीक स्पर्श बिंदु $x = 1$ पर।

\textbf{16.} नति-परिवर्तन पर \emph{दैनिक} रोगी-गिनती (\omterm{def:g12:deriv:derivative}{अवकलज}) अपनी चोटी छूती
है: वृद्धि का त्वरण रुक जाता है और मंदन शुरू हो जाता है --- यही पहला गणितीय
संकेत है कि लहर मुड़ रही है, और यह गिनतियों के गिरने से बहुत पहले आ जाता है।
उससे पहले $f'' > 0$ (हर दिन पिछले से बुरा); उसके बाद $f'' < 0$ (बढ़ती तो है, पर धीमी
पड़ती हुई)।

\textbf{17.} $S'(r) = 4\pi r - \frac{2V}{r^2} = 0$ से $r^3 = \frac{V}{2\pi} = \frac{330}{2\pi}$ मिलता है: $r \approx 3.74$ cm, और $h = \frac{330}{\pi r^2} \approx 7.49$ cm।

\textbf{18.} इष्टतम पर $h = \frac{V}{\pi r^2} = \frac{2\pi r^3}{\pi r^2}
= 2r$: ऊँचाई व्यास के बराबर, आयतन चाहे जो हो।

\textbf{19.} $S''(r) = 4\pi + \frac{4V}{r^3} > 0$: \omterm{def:g12:deriv:convex}{उत्तल}, इसलिए क्रांतिक त्रिज्या वैश्विक \omterm{def:g10:functions:extrema}{न्यूनतम} है। असली
डिब्बे पकड़, चिनाई, अलमारी में दिखावट और ऊपर-नीचे की मोटी धातु के कारण अधिक
ऊँचे बनाए जाते हैं --- इष्टतमीकरण सदा ठीक वही \omterm{def:g10:functions:extrema}{न्यूनतम} करता है जो आपने लिखा,
न कि वह जो आपके मन में था।

\textbf{20.} जीवनरक्षक: प्रतिरूप $T(x)$, अवकलन, स्नेल, \omterm{def:g12:deriv:convex}{उत्तलता}, एक सेकंड की
बचत। डिब्बा: प्रतिरूप $S(r)$, अवकलन, $h = 2r$, \omterm{def:g12:deriv:convex}{उत्तलता}, धातु की बचत। शहतीर
(\cref{pb:g11:deriv:1}): प्रतिरूप $S(w)$, अवकलन, $d = w\sqrt2$, सारणी, बढ़ी हुई दृढ़ता। और फ़र्मा के
सिद्धांत से प्रकाश स्वयं हर सतह पर यही प्रक्रिया चलाता है --- प्रकृति ही पहली
इष्टतमकर्ता थी।
\end{solution}
